\documentclass{article}

% NeurIPS convention uses numbered in-text citations ([1], [2], ...),
% not author-year. Confirmed directly from neurips_2026.sty's own
% source: it unconditionally does `\RequirePackage{natbib}` with NO
% options as part of processing \usepackage{neurips_2026} below,
% which locks natbib to its own default (author-year) the moment
% that line runs. A later \usepackage[numbers]{natbib} is too late --
% natbib is already loaded, so the option is silently dropped. The
% style file's own comments document \PassOptionsToPackage as the
% correct mechanism for this exact situation: it queues the option
% so it applies automatically whenever natbib actually loads,
% regardless of who triggers that load. Must still appear before
% \usepackage{neurips_2026} below, since that's where the load happens.
\PassOptionsToPackage{numbers}{natbib}

% NOTE TO AUTHOR: this file assumes the official neurips_2026.sty package,
% which is NOT reproduced here (it is proprietary conference formatting
% infrastructure, not paper content, and should be downloaded fresh from
% https://neurips.cc/Conferences/2026/CallForPapers to guarantee it matches
% the current official version). Place neurips_2026.sty in the same
% directory as this file before compiling. Swap the option below per your
% submission stage: no option / [preprint] for arXiv, [final] for camera-ready.
\usepackage{neurips_2026}

\usepackage[utf8]{inputenc}
\usepackage[T1]{fontenc}
\usepackage{hyperref}
\usepackage{url}
\usepackage{booktabs}
\usepackage{amsfonts}
\usepackage{amsmath}
\usepackage{amssymb}
\usepackage{amsthm}
\usepackage{nicefrac}
\usepackage{microtype}
\usepackage{xcolor}
\usepackage{graphicx}
\usepackage{subcaption}
\usepackage{multirow}
\usepackage{placeins}
\usepackage{array}

\newtheorem{lemma}{Lemma}
\newtheorem{conjecture}{Conjecture}
\newtheorem{corollary}{Corollary}
\newtheorem{definition}{Definition}
\newtheorem{remark}{Remark}

\title{The Isotypic Bottleneck: A Representation-Theoretic Bound on Rank Collapse During Grokking, and Its Limits as a Precursor Signal}

\author{%
  Md Muntaqim Meherab \\
  Department of Computer Science and Engineering \\
  Daffodil International University \\
  Dhaka, Bangladesh \\
  \texttt{[email address]}
}

% Confirmed from neurips_2026.sty's own source: it never loads hyperref
% itself, so hyperref here comes only from the explicit \usepackage{hyperref}
% line above -- placement relative to neurips_2026 doesn't matter for this
% one. colorlinks=true removes the border boxes hyperref draws
% by default around \cite/\ref/\S\ref/URLs; black keeps the printed page
% looking like plain text while the links stay live/clickable in the PDF.
\hypersetup{
  colorlinks=true,
  linkcolor=black,
  citecolor=black,
  urlcolor=black,
}

\begin{document}

\maketitle

\begin{abstract}
Effective-rank collapse has been proposed as an early, computationally cheap
precursor to grokking, but we find a puzzling asymmetry: the
signal reliably precedes generalization on modular arithmetic, yet fails on
sparse parity. We give a representation-theoretic account of this asymmetry.
We prove that if a group-equivariant representation's image lies in the
isotypic components of a set of irreducible representations, its rank is
bounded by the summed dimension of those components; for cyclic groups this
bound is trivially proportional to the number of active Fourier frequencies,
but for non-abelian groups the bound depends on \emph{which} irreducible
representations are active, not merely how many. Across three independent
GPU-verified
experimental runs on modular addition ($p \in \{61,97,113\}$, $n=15$), the
plateau rank tracks this bound closely ($r=0.87$, $p<10^{-4}$); extending the
test to a non-abelian group of matched training dynamics (the dihedral group
$D_4$) and pooling with modular arithmetic, while correctly excluding the one
family that contributes no within-family evidence, gives $r=0.97$
($p<10^{-12}$, $n=20$). We then test the natural proposal that a
representation-theoretic alignment score, rather than raw rank, restores a
usable precursor on sparse parity, and find that it does not: across seven
signal/task combinations and two independent statistical tests, no
Bonferroni-corrected result supports it. In its place we report an unrelated,
independently-motivated finding: a weight-decay-driven ``clock'' derived from
the optimizer's own decay dynamics acts as a valid, if loose, lower bound on
generalization time in $25$ of $30$ runs, and the five exceptions are exactly
the five runs that independently fail a delayed-generalization diagnostic.
Finally, we report a controlled experiment
disentangling why the symmetric group $S_4$ never converges within budget:
neither $S_4$ nor $D_4$ reaches our grokking criterion at any of five tested
weight-decay values, ruling out weight decay alone as the explanation, though
$D_4$ shows substantial partial generalization ($11$--$46\%$ validation
accuracy across ten runs, against a $12.5\%$ chance rate) where $S_4$ shows
essentially none ($3$--$9\%$ across fifteen runs, against a $4.2\%$ chance
rate). We report all of this, including the negative results, because we
believe an honest account of what a representation-theoretic bound does and
does not predict is more useful than a narrower story that only reports the
parts that worked.
\end{abstract}

\section{Introduction}
\label{sec:intro}

Grokking --- the phenomenon in which a network fits its training set almost
immediately, then generalizes only much later, after a delay that can span
orders of magnitude more optimization steps --- was first documented on small
algorithmic tasks by \citet{power2022grokking} and has since become a
tractable testbed for understanding how generalizing circuits form inside
neural networks
\citep{liu2023omnigrok,thilak2022slingshot,nanda2023progress,zhong2023clock}.
On modular addition specifically, \citet{nanda2023progress} showed that the
generalizing circuit is legible: the network learns a small set of Fourier
frequencies and composes them via trigonometric identities, and the
\emph{effective rank} of its hidden representation collapses sharply as this
circuit consolidates, well before validation accuracy visibly moves. This
collapse has been proposed as a cheap, architecture-agnostic early-warning
signal for grokking, since effective rank is far cheaper to compute during
training than held-out validation accuracy.

That proposal runs into a genuine asymmetry. On $k$-sparse parity ---
predicting the sign of a product of $k$ input bits out of $n$, a task whose
sample complexity is governed by the statistical-query lower bounds of
\citet{barak2022hidden} rather than by circuit legibility --- rank collapse
does not reliably precede generalization at all: participation ratio stays
flat across training, with no collapse before, at, or after the grokking
transition (Figure~\ref{fig:asymmetry}). This is a genuine anomaly, not a minor
caveat: the cheapest, most architecture-agnostic precursor to grokking known
works on one canonical algorithmic family and silently fails on another, with
no account of \emph{why}.

We give a representation-theoretic account of this asymmetry, and we test it
honestly, including the parts where the natural next hypothesis turned out to
be wrong. The starting observation is simple linear algebra with a
group-theoretic label attached: if a network's internal representation of a
group's elements is equivariant, its image decomposes into isotypic
components indexed by the group's irreducible representations (irreps), and
its rank is bounded by the summed dimension of whichever isotypic components
are actually occupied (\S\ref{sec:theory}, Lemma~\ref{lem:bound}). For a
cyclic group $\mathbb{Z}_p$, every non-trivial real irrep is
two-dimensional, so this bound reduces to twice the number of active Fourier
frequencies --- recovering Nanda et al.'s picture as a special case: with
only one irrep dimension available, ``how compressed the circuit is'' and
``how many frequencies it uses'' are the same question. For $k$-sparse
parity, the relevant symmetry is the hypercube sign group $(\mathbb{Z}/2)^n$,
whose irreps (Walsh characters) are all one-dimensional; the SQ-hardness
results of \citet{barak2022hidden} suggest gradient descent does not find a
correspondingly compressed solution, so no comparable collapse should be
expected, and none is observed.

If the bound is real, it should generalize beyond cyclic groups in a way
that is a genuinely stronger test than ``more primes of the same
family'': every non-trivial irrep of $\mathbb{Z}_p$ has the same
dimension, so a bound stated in terms of \emph{how many} irreps are
active is indistinguishable, for that family, from one stated in terms of
\emph{total occupied dimension}. Non-abelian groups break this
degeneracy, since irrep dimensions there are not all equal; we test $S_4$
(dimensions $1,1,2,3,3$) and $D_4$ (dimensions $1,1,1,1,2$) for exactly
this reason. This choice mattered in a way we did not anticipate: $S_4$
never converges within budget at any tested weight decay
(\S\ref{sec:disentangle}), limiting what it can contribute to the test,
which we report directly rather than work around (Figure~\ref{fig:concept},
Appendix~\ref{app:extrafigs}, illustrates the contrast this choice
exploits).

We also test, and report honestly failing, the most natural fix: since
parity's relevant structure lives in one-dimensional Walsh characters rather
than compressible higher-dimensional irreps, perhaps a signal built directly
from low-degree Walsh alignment --- rather than from rank --- would recover
a usable precursor. Across seven signal/task combinations and two
independent statistical tests with Bonferroni correction, it does not
(\S\ref{sec:h2}). In its place, we report an unrelated, independently
motivated finding that does hold up: a weight-decay-driven timing bound,
built entirely from the optimizer's own decay dynamics, functions as a
genuine lower bound on generalization time in $25$ of our $30$
modular-arithmetic and parity runs, with the specific pattern of exceptions
not an artifact of the bound's single free parameter (\S\ref{sec:clock}).

\paragraph{Contributions.}
\begin{itemize}
  \item A proven lemma bounding the rank of an isotypic-supported
  representation by the summed dimension of its occupied components
  (Lemma~\ref{lem:bound}), with a clearly labeled, unproven conjecture
  (Conjecture~\ref{conj:parity}) for why sparse parity should not satisfy
  the collapse it would otherwise predict.
  \item Confirmation of the bound's structural prediction on modular
  arithmetic, replicated across three independent full runs ($r=0.87$,
  $p<10^{-4}$, $n=15$), and its extension to a non-abelian group with a
  genuinely different irrep-dimension profile, with the pooled evidence
  reported alongside an explicit accounting of which part is load-bearing
  (\S\ref{sec:h1_results}).
  \item An honestly negative result: the natural fix for the parity
  precursor gap does not restore a usable signal, tested with two
  independent statistical procedures and full multiple-comparison
  correction (\S\ref{sec:h2}).
  \item An independent, unfitted weight-decay clock holding as a genuine
  lower bound on generalization timing where the representation-theoretic
  account is silent, with a robustness check confirming this is not an
  artifact of the bound's one free parameter (\S\ref{sec:clock}).
  \item A controlled experiment isolating whether group size or weight
  decay explains $S_4$'s failure to converge, with the result reported as
  genuinely mixed rather than resolved in either direction
  (\S\ref{sec:disentangle}).
\end{itemize}

\FloatBarrier
\section{Related Work}
\label{sec:related}

\paragraph{Grokking.} Delayed generalization on small algorithmic tasks was
introduced by \citet{power2022grokking}. \citet{liu2023omnigrok} locate the
phenomenon in a training/test loss-landscape mismatch as a function of
weight norm; \citet{thilak2022slingshot} identify a related late-training
instability in adaptive optimizers. We ask a narrower question than either:
not why grokking happens, but why one proposed early-warning signal for it
works on some tasks and not others.

\paragraph{Mechanistic accounts on group-structured tasks.}
\citet{nanda2023progress} reverse-engineer the modular-addition circuit as a
Fourier-frequency, trigonometric-identity computation; \citet{zhong2023clock}
show the same task admits qualitatively different circuits (``Clock'' and
``Pizza'') depending on hyperparameters. \citet{chughtai2023toy} generalize
the representation-theoretic account to arbitrary finite groups, showing
converged circuits organize around a small number of irreducible
representations; Lemma~\ref{lem:bound} formalizes the rank consequence of
that picture and, to our knowledge, is the first to test it as a
quantitative, dimension-weighted bound rather than a qualitative
description. \citet{stander2024grokking} reverse-engineer $S_5$/$S_6$
networks and find circuits organized around \emph{coset} structure rather
than irreducible representations directly --- a different lens we do not
reconcile with ours here (\S\ref{sec:discussion}). \citet{varma2023explaining}
offer a circuit-efficiency account of grokking timing closer in spirit to
our clock (\S\ref{sec:clock}) than to our structural bound.

\paragraph{Sparse parity and SQ hardness.} \citet{barak2022hidden} show SGD
learns $k$-sparse parity near the statistical-query information limit
\citep{kearns1998efficient}, with ``hidden progress'' invisible in the loss
but present in weight statistics before generalization. Our
Conjecture~\ref{conj:parity} restates this hardness as a reason rank
collapse specifically should not occur; we do not prove it, and
\S\ref{sec:h1_results}--\S\ref{sec:ablation} show the natural fix it
motivates does not work either. \citet{odonnell2014analysis} is the
standard reference for the Walsh-Fourier analysis our alignment score is
built on.

\paragraph{Concurrent 2026 work.} \citet{sivasankar2026circuit}
independently extends a representation-theoretic precursor account to
$S_5$, generalizing \citet{chughtai2023toy}'s Cayley-table setting; it
remains confined to group composition and does not test parity or a
dimension-heterogeneous group pairing as we do.
\citet{truong2026weightnorm,truong2026normseparation} independently
establish a \emph{calibrated} weight-decay timing law; our clock
(\S\ref{sec:clock}) uses the same decay mechanism but deliberately tests
only its unfitted, conservative lower-bound form, since we lack their
calibration data.

\paragraph{Effective rank, equivariance, architecture.} We use the
participation-ratio form of \citet{roy2007effective}'s effective rank
throughout for its noise stability (\S\ref{sec:setup}). The broader
motivation for treating symmetry as a first-class object connects to
equivariant and geometric deep learning
\citep{cohen2016group,bronstein2021geometric}, though our setting ---
small transformers that discover equivariant structure through training
rather than architectures built to enforce it --- sits closer to
mechanistic interpretability. Our architecture follows
\citet{vaswani2017attention}; all models train under the decoupled weight
decay of \citet{loshchilov2019decoupled}, the same mechanism
\S\ref{sec:clock}'s clock is derived from.

\FloatBarrier
\section{Problem Setup}
\label{sec:setup}

We study two families of algorithmic classification tasks, following the
grokking literature's standard small-transformer setup.

\textbf{Modular addition.} Given $a, b \in \mathbb{Z}_p$, predict
$(a+b) \bmod p$ from the input sequence $[a, b, \texttt{=}]$. We use
$p \in \{61, 97, 113\}$.

\textbf{Sparse parity.} Given $x \in \{-1,+1\}^n$ and a fixed hidden subset
$S \subset [n]$ with $|S|=k$, predict $\mathrm{sign}\big(\prod_{i\in S} x_i\big)$.
We use $n=40$ with $k \in \{2,3,4\}$.

Both are trained with a one-layer transformer (self-attention followed by an
MLP block, following \citealp{vaswani2017attention}) under AdamW
\citep{loshchilov2019decoupled}; sparse parity additionally uses a plain MLP
readout width of $200$. Full hyperparameters are in
Table~\ref{tab:fullconfig} (Appendix~\ref{app:fullconfig}).
A run is judged to have \emph{grokked} at step $T^*_{\mathrm{grok}}$, the
first evaluation at which validation accuracy exceeds $0.95$ and remains
above that threshold for five consecutive evaluations. The asymmetry this
paper explains --- shown directly as trained trajectories from our own
runs in Figure~\ref{fig:asymmetry}
(Appendix~\ref{app:extrafigs}) --- is exactly Lemma~\ref{lem:bound} and
Conjecture~\ref{conj:parity}'s target.

\textbf{Effective rank.} At each evaluation step we record the residual-
stream activations for all inputs and compute their participation ratio,
$\mathrm{PR}(\Phi) = \big(\sum_i \sigma_i^2\big)^2 / \sum_i \sigma_i^4$,
where $\sigma_i$ are the singular values of the activation matrix $\Phi$.
We use this form of effective rank \citep{roy2007effective} rather than
its entropy-based cousin or a hard numerical-rank threshold because a
direct noise-injection check confirms it degrades far more gracefully
under realistic noise (Appendix~\ref{app:metric_robustness}).


\FloatBarrier
\section{Theory}
\label{sec:theory}

We work with real irreducible representations throughout, following
\citet{serre1977linear}. Let $G$ be a finite group with real irreducible
representations $\rho_1,\ldots,\rho_r$ of dimensions $d_1,\ldots,d_r$.

\begin{lemma}[Isotypic rank bound]
\label{lem:bound}
Let $\Phi \in \mathbb{R}^{|G|\times D}$ be an embedding matrix, rows indexed
by $g \in G$, and suppose that $\Phi$'s column space, viewed as a
representation of $G$ under left translation of the rows, is isomorphic to
a subrepresentation of $\bigoplus_i \rho_i^{\oplus n_i}$ for multiplicities
$n_i \ge 0$. Then $\mathrm{rank}(\Phi) \le \sum_i n_i d_i$.
\end{lemma}
\begin{proof}
By hypothesis the column space of $\Phi$ is isomorphic, as an abstract
vector space, to a subspace of $\bigoplus_i \rho_i^{\oplus n_i}$, which has
real dimension $\sum_i n_i d_i$; the dimension of a subspace cannot exceed
the dimension of its ambient space. The rank of a matrix is the dimension
of its column space, so $\mathrm{rank}(\Phi) = \dim(\mathrm{col}(\Phi)) \le
\sum_i n_i d_i$.
\end{proof}

Lemma~\ref{lem:bound} is a direct consequence of dimension counting once the
isotypic hypothesis is granted; its content is in connecting rank to
representation type at all, not in the proof itself. What the bound predicts
in practice depends entirely on the multiplicities $n_i$, and here abelian
and non-abelian groups behave differently in a way we verify numerically
(Appendix~\ref{app:multiplicity}) before relying on it.

\begin{remark}[Multiplicity: forced for abelian groups, an assumption for
non-abelian ones]
\label{rem:multiplicity}
For abelian $G$, every irrep has $d_i\in\{1,2\}$ and a non-trivial one acts
as scalar multiplication, so a second independent embedding of the same
irrep provably adds no rank (verified for $\mathbb{Z}_7$,
Appendix~\ref{app:multiplicity}): $n_i=1$ is \emph{forced}, not assumed,
and Lemma~\ref{lem:bound} reduces to the familiar
$\mathrm{rank}(\Phi)\le 2m$ for $m$ active Fourier frequencies, recovering
\citet{nanda2023progress} as a special case. For non-abelian $G$ with an
irrep of dimension $d_i>1$ (e.g.\ $S_4$'s two $3$-dimensional irreps,
$D_4$'s one $2$-dimensional irrep), this argument breaks: $\rho_i(g)$ is a
genuine matrix, not a scalar, and a second independent embedding
\emph{can} add rank, up to the ceiling $d_i^2$ set by that irrep's
multiplicity in the regular representation (verified for $S_4$'s standard
representation, Appendix~\ref{app:multiplicity}: one embedding gives rank
$3$ as predicted under $n_i=1$, a second pushes it to $6$, not $3$). The
bound we test empirically for $S_4$ and $D_4$ in \S\ref{sec:h1_results}
therefore rests on an explicit, stated assumption --- multiplicity-one,
the minimal non-redundant case, motivated by but not implied by prior
findings that grokked circuits converge to minimal bases
\citep{nanda2023progress,chughtai2023toy} --- which we flag directly
rather than let the lemma's generality quietly cover a claim it does not,
on its own, make.
\end{remark}

\begin{conjecture}[Sparse parity does not satisfy the collapse]
\label{conj:parity}
For $k$-sparse parity on $n$ bits under the hypercube sign group
$(\mathbb{Z}/2)^n$ --- whose only irreducible representations are the
one-dimensional Walsh characters $\chi_S(x) = \prod_{i \in S} x_i$ --- we do
not expect gradient descent under weight decay to converge to a
representation whose column space is supported on a small subset of Walsh
characters, because \citet{barak2022hidden} show that any statistical-query
algorithm for $k$-sparse parity requires either exponentially many queries
or a query tolerance exponentially small in $n$, and gradient descent with a
bounded number of directions behaves as an SQ algorithm in the relevant
regime. If the SQ-hardness result forces the optimizer to maintain at least
$k$ approximately orthogonal directions throughout training rather than
consolidating onto the target's own low-dimensional spectral support, no
rank collapse of the kind predicted for finite groups by
Lemma~\ref{lem:bound} under Remark~\ref{rem:multiplicity}'s abelian-case
reasoning should occur, and the representation's rank should remain high
throughout training. This is a hypothesis, not a theorem: we do not prove
that SQ hardness forces this outcome, only that it is a plausible
mechanism consistent with the failure we document empirically
(Figure~\ref{fig:asymmetry}), and we test it further alongside its most
natural fix (\S\ref{sec:h1_results}--\S\ref{sec:h2}).
\end{conjecture}

\FloatBarrier
\subsection{The weight-decay clock}
\label{sec:clock_theory}

Independently of the representation-theoretic account above, we test a
second, unrelated timing mechanism motivated by
\citet{truong2026weightnorm,truong2026normseparation}: under AdamW's
decoupled weight decay \citep{loshchilov2019decoupled}, a parameter norm
shrinks multiplicatively by a factor $(1-\eta\lambda)$ per step in the
absence of a compensating gradient signal, giving $\approx
e^{-\eta\lambda t}$ decay after $t$ steps. The number of steps for such a
quantity to decay to a fraction $\rho$ of its initial value is
\begin{equation}
\tau_c = \frac{-\ln \rho}{\eta \lambda}.
\label{eq:clock}
\end{equation}
We deliberately test only this bare, unfitted form, with a fixed, unfit
$\rho=0.8$, rather than the calibrated multiplier $\bar{c}$ that
\citet{truong2026weightnorm} establish on top of it: we do not have that
calibration procedure or its fitting data, and we are not willing to
reconstruct it. What we test in \S\ref{sec:clock} is the weaker, honest
claim that $\tau_c$ is a valid \emph{lower bound} on generalization time,
$T^*_{\mathrm{grok}} > \tau_c$, not that $T^*_{\mathrm{grok}} \approx
\bar{c}\,\tau_c$ for any particular fitted $\bar{c}$.

\section{Experimental Setup}
\label{sec:expsetup}

Complete hyperparameters for every task and every sweep are in
Appendix~\ref{app:fullconfig}, Table~\ref{tab:fullconfig}, which we
cross-checked field-by-field against the exact \texttt{config.json} logged
by the run; we summarize only what is needed to read the results below.
All tasks use AdamW \citep{loshchilov2019decoupled}, $\eta=10^{-3}$,
$\beta=(0.9,0.98)$, and $5$ seeds unless noted; weight decay is $2.0$
(modular, $S_4$, $D_4$) or $0.7$ (parity), except where explicitly swept
(\S\ref{sec:disentangle}, \S\ref{sec:robustness}).

\textbf{Architecture.} Modular addition and both group-composition tasks use
a one-layer, four-head decoder-only transformer
\citep{vaswani2017attention} with $d_{\mathrm{model}}=128$ and an MLP width
of $512$; sparse parity uses a plain two-layer MLP of width $200$, following
\citet{barak2022hidden}'s setup for SQ-hardness demonstrations rather than a
transformer, since the task has no sequential structure to attend over.
Parameter counts, computed directly from the instantiated model classes and
cross-checked against the values logged during training, are $214{,}656$
(modular, largest prime), $8{,}602$ (parity), $205{,}184$ ($S_4$), and
$201{,}088$ ($D_4$).

\textbf{Metrics.} Effective rank is the participation ratio of
\S\ref{sec:setup}. Fourier concentration (modular arithmetic) is the
fraction of DFT power in the top-$4$ non-DC frequencies of the token
embedding. Isotypic energy fractions ($S_4$, $D_4$) come from the
character-theoretic projection $\Pi_\rho = \frac{d_\rho}{|G|}\sum_{g}
\chi_\rho(g)\,L_g$ on the token embedding ($L_g$: the $|G|\times|G|$
permutation matrix for left multiplication by $g$) --- the direct
isotypic decomposition \S\ref{sec:theory} is about, not a proxy for it.
The predicted bound is the summed dimension of the fewest, highest-energy
irreps reaching $90\%$ of total energy; for $\mathbb{Z}_p$ we use the
dimension-comparable analogue ($2$ dimensions per DFT frequency, same
threshold), putting both families on the same predicted-dimension axis
for the pooled test in \S\ref{sec:h1_results}.

The Walsh alignment score (sparse parity) fits a degree-bounded Walsh
basis to hidden activations and reports out-of-sample $R^2$ (half a
dedicated probe batch to fit, the held-out half to score, decoupled from
the task's own validation set). This is necessary, not a convenience: the
in-sample version substantially overfits pure noise at our probe sizes
(Appendix~\ref{app:walsh_bias}); every Walsh-score number reported in
this paper uses only the out-of-sample construction.

\textbf{Statistical protocol.} For each candidate precursor signal we
compute a lead time $\Delta T = T^*_{\mathrm{grok}} - T^*_{\mathrm{signal}}$
per run, where $T^*_{\mathrm{signal}}$ is the first step at which the signal
crosses the midpoint of its own $10^{\mathrm{th}}$--$90^{\mathrm{th}}$
percentile range, with the crossing direction determined by comparing
early- and late-window means. We test $\Delta T > 0$ with a one-sided
Wilcoxon signed-rank test, cross-checked with a sign test, report
rank-biserial correlation as an effect size, and apply Bonferroni
correction across all comparisons reported together ($m=7$, $\alpha =
0.05/7 \approx 0.00714$).

\textbf{Compute.} All experiments run on a single NVIDIA Tesla T4 (peak
memory $780$\,MB); the $70$ training runs underlying this paper total
$15{,}205$ wall-clock seconds ($4.2$ hours), well within free-tier GPU
access.


\section{Results}
\label{sec:results}


\subsection{H1: rank tracks the isotypic bound}
\label{sec:h1_results}

On modular arithmetic, plateau participation ratio correlates strongly with
the predicted bound: $r=0.87$ ($p=3.0\times10^{-5}$, $n=15$), replicated
identically across three independent full GPU runs. Extending the axis to
$S_4$ and $D_4$ and pooling all three families gives $r=0.73$
($p=4.0\times10^{-5}$, $n=25$) --- technically significant, but this number
should never be quoted alone. $S_4$ contributes no within-family test:
its predicted bound is the same constant ($8$) across all five seeds,
since $S_4$ never converges (\S\ref{sec:disentangle}), so its embedding
never moves past a generic, dimension-proportional energy distribution.
A cluster with zero within-family variance can only pull a
pooled fit as a fixed point; it cannot be tested. Excluding it and pooling
only $\mathbb{Z}_p$ with $D_4$ --- the two families that actually converge
and show genuine within-family spread --- gives $r=0.97$
($p=4.3\times10^{-13}$, $n=20$), the strongest and most defensible
cross-family number we report. $D_4$'s own within-family correlation is
directionally positive but not significant on its own ($r=0.63$, $p=0.25$,
$n=5$): five points is not enough statistical power to confirm a
relationship on that family alone, only to be consistent with one.
Table~\ref{tab:h1honesty} reports all four numbers together, which is how
we believe this result should be read and cited. An alternative,
concentration-based version of the $\mathbb{Z}_p$-only test (Fourier
concentration rather than the dimension-weighted bound) gives the same
picture even more cleanly ($r=-0.97$, $p=1.7\times10^{-9}$, $n=15$;
Figure~\ref{fig:h1zp}, Appendix~\ref{app:extrafigs}).

\begin{table}[t]
\caption{H1, reported with the full breakdown a bare pooled correlation
would hide. The pooled number is real but is substantially carried by
between-cluster separation rather than a single relationship holding
smoothly across the full range; the $\mathbb{Z}_p$-only and
excl.-$S_4$ rows are the load-bearing evidence.}
\label{tab:h1honesty}
\centering
\small
\begin{tabular}{@{}lrrr@{}}
\toprule
Comparison & $r$ & $p$ & $n$ \\
\midrule
Pooled ($\mathbb{Z}_p$ + $S_4$ + $D_4$) & $0.73$ & $4.0\times10^{-5}$ & $25$ \\
Pooled excluding $S_4$ ($\mathbb{Z}_p$ + $D_4$) & $\mathbf{0.97}$ & $4.3\times10^{-13}$ & $20$ \\
$\mathbb{Z}_p$ only (load-bearing) & $0.87$ & $3.0\times10^{-5}$ & $15$ \\
$D_4$ only & $0.63$ & $0.25$ & $5$ \\
$S_4$ only & \multicolumn{3}{c}{undefined (zero within-family variance)} \\
\bottomrule
\end{tabular}
\end{table}

\begin{figure}[t]
\centering
\includegraphics[width=0.95\textwidth]{figures/fig5_h1_nonabelian_generalization.pdf}
\caption{H1 across group families. \textbf{Left:} predicted isotypic bound
vs.\ measured plateau rank for $\mathbb{Z}_p$ (circles), $S_4$ (squares),
and $D_4$ (triangles), with the pooled and $\mathbb{Z}_p$-only correlations
annotated together on the plot rather than leaving the pooled number to
stand alone. \textbf{Right:} the isotypic energy fraction $S_4$'s embedding
places on each of its five irreps at the plateau, mean $\pm$ one standard
deviation across seeds --- the concrete complement to the left panel's
summary statistic, and the pattern (energy concentrated on the
$3$-dimensional irreps regardless of seed) consistent with $S_4$ reflecting
a generic, dimension-proportional prior rather than task-driven
specialization, as expected from a network that never converges.}
\label{fig:h1}
\end{figure}

\subsection{H2: the representation-theoretic fix does not restore a parity precursor}
\label{sec:h2}

Table~\ref{tab:h2} reports all seven signal/task comparisons. The
headline test --- whether the out-of-sample Walsh alignment score
(degree $\le 1$) leads grokking on sparse parity --- is not supported:
mean lead $573$ steps, $95\%$ CI $[-1507, 3307]$ (crosses zero), Wilcoxon
$p=0.744$, only $5$ of $15$ runs positive, not Bonferroni-significant. The
degree-$2$ Walsh score is not supported either, and by a wider margin
($p=0.819$, only $4/15$ positive, $r=-0.27$). The one comparison where the
two statistical tests disagree is weight norm on parity: the Wilcoxon test
nominally crosses the Bonferroni threshold ($p=0.005$), but the sign test
does not ($p=0.151$, only $10/15$ positive) --- exactly the kind of
disagreement a robust cross-check is meant to catch, and we read it as
Wilcoxon's sensitivity to the magnitude of a few large positive differences
rather than a consistent direction across most runs, not as a confirmed
result. Rank does not lead on parity either ($r=-0.57$). Weight norm is the
only signal with any positive lead evidence on parity at all ($p=0.005$,
$10/15$ positive, $r=0.75$, with the disagreement above noted), and it is
the mechanism already established by
\citet{truong2026weightnorm,truong2026normseparation}, not this paper's
proposed contribution.

\begin{table}[t]
\caption{Lead-time significance, all seven signal/task comparisons, one-sided
Wilcoxon signed-rank test against zero with sign-test cross-check and
Bonferroni correction ($\alpha=0.05/7$). Positive mean lead indicates the
signal crosses its threshold before grokking. $\dagger$: nominally crosses
the Bonferroni threshold but disagrees with its own sign-test cross-check
($p=0.151$); read as Wilcoxon's sensitivity to a few large-magnitude
differences rather than a consistent direction, not a confirmed result
(\S\ref{sec:h2}).}
\label{tab:h2}
\centering
\footnotesize
\begin{tabular}{@{}llrrrrc@{}}
\toprule
Task & Signal & $n_+/n$ & Wilcoxon $p$ & Sign $p$ & $r$ & Sig.\ \\
\midrule
Modular & participation ratio & $11/15$ & $0.015$ & $0.059$ & $0.67$ & \ding{55} \\
Modular & weight norm & $7/15$ & $0.098$ & $0.696$ & $0.44$ & \ding{55} \\
Modular & Fourier concentration & $0/15$ & $0.9997$ & $1.000$ & $-1.00$ & \ding{55} \\
Parity & participation ratio & $2/15$ & $0.971$ & $0.9995$ & $-0.57$ & \ding{55} \\
Parity & weight norm & $10/15$ & $\mathbf{0.005}^\dagger$ & $0.151$ & $0.75$ & (\checkmark) \\
Parity & Walsh score (deg.\ $\le\!2$) & $4/15$ & $0.819$ & $0.982$ & $-0.27$ & \ding{55} \\
Parity & Walsh score (deg.\ $\le\!1$, primary) & $5/15$ & $0.744$ & $0.941$ & $-0.19$ & \ding{55} \\
\bottomrule
\end{tabular}
\end{table}

\begin{figure}[t]
\centering
\includegraphics[width=0.95\textwidth]{figures/fig3_h2_evidence.pdf}
\caption{H2 evidence on sparse parity. \textbf{Left:} the Walsh alignment
score (degree $\le 1$) trajectory, mean $\pm 95\%$ CI, with the mean
grokking onset marked. \textbf{Middle:} lead-time distributions by signal
for both task families, with the primary comparison's effect size and
Bonferroni-corrected significance annotated on the plot. \textbf{Right:}
the same parity runs' rank lead and Walsh lead, paired and connected per
run --- the aggregate in the middle panel can hide whether an individual
run's two signals actually agree in direction; this panel shows that
directly rather than assuming the aggregate speaks for every run.}
\label{fig:h2}
\end{figure}

\subsection{The $k$-sweep breaks the ``parity never collapses'' picture in a specific way}
\label{sec:ablation}

Averaging H2 over $k\in\{2,3,4\}$ hides a pattern that matters for how
Conjecture~\ref{conj:parity} should be read (full per-$k$ table,
Appendix~\ref{app:extratables}). At $k=2$, plateau rank collapses to
$2.85$ --- lower than several of the modular-arithmetic plateaus in
Table~\ref{tab:h1honesty} --- and the Walsh-score lead is strongly negative
($-640$ steps). We do not read this as parity failing Conjecture~\ref{conj:parity};
we read it as $k=2$ failing the conjecture's own precondition. All five
$k=2$ seeds reach our grokking threshold at $t=200$ steps --- the first
evaluation point our logging captures at all, and identical across every
seed --- which is the signature of a task solved immediately, not of a
memorization-then-generalization delay. SQ-hardness arguments that predict
no rank collapse under a genuine delayed-generalization regime say nothing
about a task that never enters that regime in the first place. $k=3$ and
$k=4$ look nothing like this (plateau rank $29.4$ and $15.2$,
$t_{\mathrm{grok}}$ ranging from $3{,}600$ to $15{,}400$ steps across
seeds), and it is exactly these two $k$ values whose clock ratios hold in
\S\ref{sec:clock} while every one of the five $k=2$ runs' does not --- two
independently derived diagnostics, rank collapse and the weight-decay
clock, converging on the same five runs as the exception, which we read
as corroboration rather than coincidence.


\subsection{The weight-decay clock: an independent, honest lower bound}
\label{sec:clock}

The clock of Eq.~\eqref{eq:clock}, evaluated at the unfitted $\rho=0.8$,
gives $\tau_c=111.6$ steps for modular arithmetic and $\tau_c=318.8$ steps
for parity, using each task's own $(\eta,\lambda)$. Across all $30$ runs
where a grokking step is defined, $T^*_{\mathrm{grok}} > \tau_c$ in $25$
--- every one of the $15$ modular-arithmetic runs (ratio range
$10.8\times$--$44.8\times$) and $10$ of the $15$ parity runs. The five
exceptions are exactly the five $k=2$ runs discussed in
\S\ref{sec:ablation} (ratio $0.63\times$, uniformly, since $T^*_{\mathrm{grok}}=200$
for all five); $k=3$ and $k=4$ hold with the same margin as modular
arithmetic (ratio range $11.3\times$--$48.3\times$). We regard this bound
as loose --- a factor of ten to fifty above the tightest possible bound is
not a tight prediction --- and that looseness is expected given we
deliberately did not fit $\bar{c}$ (\S\ref{sec:clock_theory}); what we take
as informative is that the bound holds cleanly on every run with a genuine
delayed-generalization transition and fails on every run without one.

\begin{figure}[t]
\centering
\includegraphics[width=0.95\textwidth]{figures/fig6_weight_decay_clock.pdf}
\caption{The weight-decay clock as a lower bound. \textbf{Left:}
$T^*_{\mathrm{grok}}$ vs.\ $\tau_c$ by family, with the $y=x$ line marking
where the bound would be tight; points sit above it, not on it, consistent
with an intentionally unfitted, conservative bound. \textbf{Right:} the
$T^*/\tau_c$ ratio distribution by family.}
\label{fig:clock}
\end{figure}

\paragraph{Robustness to the bound's one free parameter.} Since $\rho=0.8$
was a stated default, not fitted, we recompute the bound across
$\rho\in\{0.1,0.3,0.5,0.8,0.9\}$ --- a full order of magnitude in $\tau_c$
(full table, Appendix~\ref{app:extratables}). Modular arithmetic holds at
$15/15$ throughout; parity holds at exactly $10/15$ for every value except
the most lenient ($\rho=0.9$), where the bound becomes too loose to be
informative. This is the behavior of a genuine effect, not an artifact of
one convenient parameter: the $k=2$ exception survives a full order of
magnitude of strictness and disappears only when the bound stops testing
anything at all.

\subsection{Disentangling $S_4$'s non-convergence}
\label{sec:disentangle}

Neither $S_4$ nor $D_4$ reaches our grokking threshold at any tested
weight decay: $D_4$ at $\{1.0, 2.0\}$ and $S_4$ at $\{0.5, 1.0, 2.0\}$, five
seeds each (full per-condition table, Appendix~\ref{app:extratables}). This
rules out weight decay alone as $S_4$'s explanation, though the result is
mixed rather than a clean resolution: $D_4$ shows substantial, consistent
partial generalization across all ten of its runs (validation accuracy
$11$--$46\%$, mean $28.5\%$, against a $12.5\%$ chance rate for
$8$-way classification), while $S_4$ shows essentially none across all
fifteen of its runs ($3$--$9\%$, mean $5.7\%$, against a $4.2\%$ chance
rate for $24$-way classification). Group size or complexity plausibly
matters --- $D_4$ is a third the elements and a ninth the total training
pairs of $S_4$ --- but partial generalization without full convergence is
not evidence for a specific mechanism. We report this as genuinely
unresolved rather than resolve it in either direction.

\subsection{Robustness: rank tracks structure, not optimizer speed}
\label{sec:robustness}

A structural claim should not depend on the optimizer knob that sets
training \emph{speed}. Sweeping modular arithmetic's smallest prime across
weight decay $\{0.5,1.0,3.0,4.0,6.0\}$ (three seeds each; full results in
Figure~\ref{fig:h1zp}, Appendix~\ref{app:extrafigs}), plateau rank among
the nine converging runs stays in a comparable range ($14.0$--$23.5$)
while $T^*_{\mathrm{grok}}$ varies far more and non-monotonically, from
$2{,}800$ to $15{,}600$ steps, with complete non-convergence at
$\lambda=6.0$. Stable structure alongside timing that varies by more than
$5\times$ and fails outright at the high end is what a representation-
theoretic account predicts and a pure-timing account would not; with only
six of fifteen sweep runs failing to converge, we treat this as
suggestive, not independent statistical weight, for \S\ref{sec:h1_results}.

\FloatBarrier
\section{Discussion and Limitations}
\label{sec:discussion}

\textbf{H2 is a genuine negative result.} We checked it three ways before
accepting it: two independent statistical tests that agree it is not
supported, a $k$-sweep showing a real $k$-dependent split, not
averaged-out noise, and replication across three independent full runs
with identical results. We have no positive account of what
\emph{does} restore a precursor on $k$-sparse parity beyond the weight-norm
mechanism prior work already established
\citep{truong2026weightnorm,truong2026normseparation}; closing that gap is
the natural next question, not one this paper answers.

\textbf{The clock is deliberately the unfitted, weaker claim.} The bound we
find is loose (ratios $10\times$--$50\times$, not near $1\times$). A
tighter, calibrated version almost certainly exists --- that is the
content of \citet{truong2026weightnorm} --- but we lack its calibration
procedure, and testing a weaker claim honestly seemed preferable to
reconstructing a stronger one we could not verify.

\textbf{Multiplicity-one, and $S_4$ vs.\ $D_4$, are open questions we
narrowed but did not close.} Remark~\ref{rem:multiplicity} and
Appendix~\ref{app:multiplicity} show a second independent embedding of a
non-abelian irrep can add rank past the bound we test; if $S_4$ or $D_4$
converges under a different recipe, the bound needs re-checking, not
assumed. The disentangling grid (\S\ref{sec:disentangle}) rules out
weight decay alone as $S_4$'s blocker; we see architecture depth/width
and longer training at $D_4$'s scale (already extended fivefold for
$S_4$, with no effect) as the two most promising next steps.

\textbf{Scale, and relation to coset-based accounts.} All experiments run
on a single T4 GPU within four hours, deliberately --- but this means our
largest task ($S_4$, order $24$) is still small next to the literature
reverse-engineering \emph{converged} $S_5$/$S_6$ circuits
\citep{stander2024grokking}. We do not reconcile our isotypic account
with their coset-based one (\S\ref{sec:related}): both may describe the
same solution differently, and since neither $S_4$ nor $D_4$ converges
here, we can only say, via Lemma~\ref{lem:bound}, what a compatible bound
would look like if they did.


\section{Conclusion}
\label{sec:conclusion}

We set out to explain why rank collapse precedes grokking reliably on
modular arithmetic but not on sparse parity. A representation-theoretic
bound (Lemma~\ref{lem:bound}) explains the asymmetry's structural half
cleanly, confirmed across three independent replications and extended,
honestly, to a non-abelian group. It does not explain the timing half:
the natural fix we proposed for parity does not work, tested two
independent ways. In its place, an unrelated, independently derived
weight-decay clock provides a genuine --- if loose --- account of
timing, with a robustness property (\S\ref{sec:clock}) we verified
rather than assumed. We think a bound honestly reported as explaining
part of what it was built to explain, not all of it, is more useful to
the next person who tests it than a narrower story would have been.


\bibliographystyle{plainnat}
\bibliography{references}

\newpage
\appendix

\section*{Appendix Overview}
\addcontentsline{toc}{section}{Appendix Overview}

The appendix is organized to mirror the paper's own argument rather than
simply collecting leftovers. Appendix~\ref{app:extrafigs} and
Appendix~\ref{app:extratables} hold figures and tables that are discussed
in the main body but whose full form did not fit within the page limit ---
every number in them is cited exactly as shown here, nothing is
summarized or approximated on the way in. Appendices~\ref{app:multiplicity}
and~\ref{app:chartables} give the mathematical machinery behind
\S\ref{sec:theory} in full: the numerical checks that justify treating
multiplicity-one as forced for abelian groups and as an assumption for
non-abelian ones, and the two character tables the entire $S_4$/$D_4$
analysis depends on, each verified two independent ways before being
trusted. Appendices~\ref{app:metric_robustness} and~\ref{app:walsh_bias}
document two methodological decisions --- which effective-rank statistic
to use, and why the Walsh alignment score needs an out-of-sample
correction --- that shaped every number reported in \S\ref{sec:results}.
Appendix~\ref{app:s4failure} gives the qualitative picture behind
\S\ref{sec:disentangle}'s headline claim, and Appendix~\ref{app:fullconfig}
closes with a single, complete reference table for every hyperparameter,
sweep, and seed used anywhere in this paper.

\FloatBarrier
\section{Additional Figures}
\label{app:extrafigs}

The three figures below are referenced from the main body but placed here
to keep the main body within the page limit; none of them are new
content, and each is discussed in full at the point the main text cites
it. We include them here, rather than describing them only in words,
because a rank-collapse claim is easiest to check by eye, and because the
qualitative pattern in each figure --- a circle of frequencies collapsing
to a handful of points, a rank trajectory falling and staying flat, a
robustness sweep that moves timing without moving structure --- carries
information a summary statistic alone does not.

Figure~\ref{fig:concept} illustrates the structural contrast
\S\ref{sec:intro} motivates the paper with: why a network solving modular
arithmetic has a small subspace available to collapse into, and why a
network solving sparse parity does not. Figure~\ref{fig:asymmetry} makes
that same contrast empirical, plotting the trained rank trajectories
themselves rather than describing them. Figure~\ref{fig:h1zp} then closes
the loop on H1 (\S\ref{sec:h1_results}): its left panel is an independent,
concentration-based version of the $\mathbb{Z}_p$-only correlation, and
its right panel is the weight-decay robustness check of
\S\ref{sec:robustness}, showing structure staying fixed while timing
moves --- the dissociation a purely representation-theoretic account
predicts and a pure-timing account does not.

\begin{figure}[t]
\centering
\includegraphics[width=0.85\textwidth]{figures/fig1_isotypic_bottleneck.pdf}
\caption{The structural asymmetry Lemma~\ref{lem:bound} formalizes.
\textbf{Left:} modular arithmetic's non-trivial irreps are all
two-dimensional rotations; a converged circuit using a handful of
frequencies is confined to a low-dimensional subspace by construction.
\textbf{Right:} sparse parity's irreps (Walsh characters) are all
one-dimensional and combinatorially numerous; there is no analogous small
subspace for an SQ-hard solution to collapse into
(Conjecture~\ref{conj:parity}).}
\label{fig:concept}
\end{figure}

This structural picture is a prediction until it is checked against
trained networks. Figure~\ref{fig:asymmetry} shows exactly that check,
using the same participation-ratio statistic used throughout
\S\ref{sec:results}, aggregated across every seed in each family.

\begin{figure}[t]
\centering
\includegraphics[width=0.95\textwidth]{figures/fig2_rank_collapse_contrast.pdf}
\caption{The documented asymmetry, mean $\pm 95\%$ CI, dotted line marking
mean grokking onset. \textbf{Left:} modular arithmetic's participation
ratio collapses as the circuit consolidates. \textbf{Right:} sparse
parity shows no comparable collapse around its own grokking transition.
This asymmetry, observed directly in our own runs, is
the phenomenon Lemma~\ref{lem:bound} and Conjecture~\ref{conj:parity}
explain.}
\label{fig:asymmetry}
\end{figure}

Finally, Figure~\ref{fig:h1zp} reports two checks on H1 that are secondary
to the pooled cross-family test of \S\ref{sec:h1_results} but strengthen
it: an alternative structural statistic that agrees with the
dimension-weighted bound, and a robustness sweep that would break the
representation-theoretic story if it failed.

\begin{figure}[t]
\centering
\includegraphics[width=0.95\textwidth]{figures/fig4_h1_and_dissociation.pdf}
\caption{H1 on $\mathbb{Z}_p$ alone, and the structure/timing dissociation.
\textbf{Left:} plateau rank vs.\ Fourier concentration (an alternative,
concentration-based structural statistic rather than the dimension-weighted
bound used elsewhere in this section), $r=-0.97$
($p=1.7\times10^{-9}$, $n=15$) --- more concentrated spectral support
predicts a smaller collapsed rank, the same relationship
Table~\ref{tab:h1honesty} tests on the dimension-weighted axis. \textbf{Right:}
the weight-decay robustness check of \S\ref{sec:robustness} --- plateau
rank stays flat across weight decay while $T^*_{\mathrm{grok}}$ shifts, the
dissociation a purely representation-theoretic account predicts and a
pure-timing account does not.}
\label{fig:h1zp}
\end{figure}

\section{Additional Results Tables}
\label{app:extratables}

Three tables are collected here, each the full version of a result
summarized in prose in the main body. In every case the summary in
\S\ref{sec:results} states the specific numbers that matter for the
paper's argument; what follows here is the complete breakdown those
summaries are drawn from, so a reader who wants to check a particular
seed, sparsity level, or weight-decay value against the aggregate claim
can do so directly.

Table~\ref{tab:ksweep} is the source of \S\ref{sec:ablation}'s central
observation: that averaging over $k$ hides a real, $k$-dependent split
rather than genuine noise, with $k=2$ behaving nothing like $k=3$ or
$k=4$ on either statistic reported.

\begin{table}[t]
\caption{Ablation over sparsity $k$ (sparse parity, $n=5$ seeds per $k$).}
\label{tab:ksweep}
\centering
\small
\begin{tabular}{@{}lrrr@{}}
\toprule
$k$ & Mean plateau rank & Mean Walsh-score (deg.\ $\le\!1$) lead & $n$ \\
\midrule
$2$ & $2.85$ & $-640$ & $5$ \\
$3$ & $29.37$ & $1960$ & $5$ \\
$4$ & $15.20$ & $-1440$ & $5$ \\
\bottomrule
\end{tabular}
\end{table}

Table~\ref{tab:rho} is the source of \S\ref{sec:clock}'s robustness claim
for the weight-decay clock: the fraction of runs clearing the bound at
each of five tested values of the clock's one free parameter $\rho$,
spanning a full order of magnitude. The pattern that matters is that
parity's shortfall from a perfect $15/15$ is the \emph{same} $10/15$ at
every value except the most lenient one tested, where the bound becomes
too loose to be informative at all.

\begin{table}[t]
\caption{Fraction of runs where $T^*_{\mathrm{grok}} > \tau_c$, by family
and $\rho$. The $k=2$ exception (parity's shortfall from $15/15$) persists
across a full order of magnitude in $\rho$ and vanishes only at the most
lenient value tested.}
\label{tab:rho}
\centering
\small
\begin{tabular}{@{}lccccc@{}}
\toprule
Family & $\rho=0.1$ & $\rho=0.3$ & $\rho=0.5$ & $\rho=0.8$ & $\rho=0.9$ \\
\midrule
Modular & $15/15$ & $15/15$ & $15/15$ & $15/15$ & $15/15$ \\
Parity & $10/15$ & $10/15$ & $10/15$ & $10/15$ & $15/15$ \\
\bottomrule
\end{tabular}
\end{table}

Table~\ref{tab:grid} is the source of \S\ref{sec:disentangle}'s
disentangling result: grokking rate and mean validation accuracy for
every group/weight-decay combination tested, showing that non-convergence
is uniform across the whole grid, while validation accuracy itself is not
--- the per-condition detail behind the $11$--$46\%$ vs.\ $3$--$9\%$
contrast reported in the main text.

\begin{table}[t]
\caption{Disentangling grid: grokking rate and mean final validation
accuracy by group and weight decay, five seeds per cell. Neither group
reaches our $0.95$ grokking threshold at any tested value.}
\label{tab:grid}
\centering
\small
\begin{tabular}{@{}lccc@{}}
\toprule
Group & Weight decay & Grok rate & Mean final val.\ acc. \\
\midrule
$D_4$ & $1.0$ & $0/5$ & $30.8\%$ \\
$D_4$ & $2.0$ & $0/5$ & $26.2\%$ \\
$S_4$ & $0.5$ & $0/5$ & $5.7\%$ \\
$S_4$ & $1.0$ & $0/5$ & $4.9\%$ \\
$S_4$ & $2.0$ & $0/5$ & $6.6\%$ \\
\bottomrule
\end{tabular}
\end{table}

\section{Multiplicity Verification (Remark~\ref{rem:multiplicity})}
\label{app:multiplicity}

Section~\ref{sec:theory} states, but does not prove in the main text, a
distinction that turns out to matter a great deal for how the results in
\S\ref{sec:h1_results} should be read: multiplicity-one is a forced
consequence of group structure for abelian groups like $\mathbb{Z}_p$, but
only an assumption --- motivated, not implied --- for non-abelian groups
like $S_4$ and $D_4$. This appendix is where that distinction is actually
checked, rather than merely asserted, because Lemma~\ref{lem:bound} is
true regardless of which case applies and it would be easy to let the
lemma's generality quietly stand in for a claim it does not by itself
make. We check both directions of the distinction directly: does a second,
independently embedded copy of the same irrep add rank, or not, for a
representative abelian and a representative non-abelian group?

\textbf{$\mathbb{Z}_7$ (abelian).} We build a $7\times 8$ embedding matrix
using frequency $k=1$'s two-dimensional real irrep, applied to one
reference vector $v_0\in\mathbb{R}^2$ and embedded into $\mathbb{R}^8$ via
a random weight matrix; its numerical rank is $2$, exactly the
multiplicity-one prediction. We then add a second, independently embedded
copy of the \emph{same} frequency --- an independent reference vector
$v_0'$ and independent embedding weights, summed into the same matrix ---
and check whether this pushes the rank past $2$. It does not: the rank
stays at $2$. The reason is structural, not a coincidence of this
particular random draw: for an abelian group, a non-trivial real irrep's
action $\rho_k(t)$ is multiplication by a single complex scalar, so
$t \mapsto \rho_k(t)v_0$ and $t \mapsto \rho_k(t)v_0'$ are both real-linear
combinations of the \emph{same} pair of basis functions of $t$ ($\cos$ and
$\sin$ at frequency $k$) regardless of which reference vector is used ---
there is no way to choose $v_0'$ that escapes this $2$-dimensional span.
This is why $n_i=1$ is forced rather than assumed for $\mathbb{Z}_p$: the
group structure itself leaves no room for a second copy to matter.

\textbf{$S_4$ (non-abelian).} We repeat the same experiment with $S_4$'s
$3$-dimensional standard representation: one reference vector
$v_0\in\mathbb{R}^3$, embedded into $\mathbb{R}^{10}$, gives a $24\times
10$ matrix of numerical rank $3$, again matching the multiplicity-one
prediction exactly. Adding a second, independently embedded copy of the
\emph{same} irrep --- an independent reference vector $v_0'$, independent
weights --- this time pushes the rank to $6$, not $3$. The mechanism is
the mirror image of the abelian case: because $S_4$ is non-abelian,
$\rho(g)$ is a genuine $3\times 3$ matrix rather than a scalar, so distinct
reference vectors are \emph{not} forced into the same $3$-dimensional
function space, and a second copy can add real rank --- up to the full
ceiling of $d_i^2=9$ set by that irrep's own multiplicity in the regular
representation, though our two-copy check here only demonstrates that the
ceiling is above $3$, not that it is reached. This is the concrete sense in
which multiplicity-one is an assumption for $S_4$ and $D_4$ rather than a
forced consequence: nothing in the group's structure alone rules out a
higher-rank solution, and the bound we test empirically in
\S\ref{sec:h1_results} depends on the network's converged solution
happening not to use one.

All ranks reported above are computed with
\texttt{numpy.linalg.matrix\_rank} at tolerance $10^{-8}$ on
double-precision matrices built from i.i.d.\ standard normal reference
vectors and embedding weights, seed $0$; we do not report a sweep over
seeds here because the underlying claim is about the achievable rank of a
specific linear-algebraic construction, not a statistical property that
varies run to run.

\FloatBarrier
\section{$S_4$ and $D_4$ Character Table Verification}
\label{app:chartables}

Every isotypic-energy computation in this paper --- the predicted bound
in Table~\ref{tab:h1honesty}, the irrep-energy breakdown in
Figure~\ref{fig:h1}, the sanity checks below --- ultimately rests on two
hand-assembled character tables, for $S_4$ and $D_4$. A character table is
a small object, easy to transcribe correctly by eye and just as easy to
get subtly wrong: a single transposed sign or a character placed under
the wrong conjugacy class would silently propagate into every downstream
number without producing an obviously invalid result. We treat this as
serious enough to check two independent ways before trusting either
table for anything else in the paper. The first check, Burnside's
identity ($\sum_i d_i^2 = |G|$), confirms the dimensions are internally
consistent but would not catch a sign error or a transposed entry, since
it only sums squared dimensions and never touches the character values
themselves. The second, full row orthogonality with conjugacy-class-size
weighting ($\sum_{c} |c|\,\chi_{\rho}(c)\overline{\chi_{\rho'}(c)} =
|G|\,\delta_{\rho\rho'}$ for every pair of irreps $\rho,\rho'$), is a much
stronger check: it must hold exactly for every pair of rows, and a single
wrong sign anywhere in the table will violate it. Both tables below pass
both checks exactly, not approximately.

\begin{table}[h]
\caption{$S_4$ character table (conjugacy classes: identity, transpositions
($6$), double transpositions ($3$), $3$-cycles ($8$), $4$-cycles ($6$)).
Verified: $\sum d_i^2 = 1+1+4+9+9=24=|S_4|$; all $\binom{5}{2}+5=15$
row-orthogonality relations hold exactly.}
\label{tab:s4chars}
\centering
\small
\begin{tabular}{@{}lrrrrr@{}}
\toprule
Irrep (dim) & $e$ & $(12)$ & $(12)(34)$ & $(123)$ & $(1234)$ \\
\midrule
Trivial ($1$) & $1$ & $1$ & $1$ & $1$ & $1$ \\
Sign ($1$) & $1$ & $-1$ & $1$ & $1$ & $-1$ \\
Two-dim.\ ($2$) & $2$ & $0$ & $2$ & $-1$ & $0$ \\
Standard ($3$) & $3$ & $1$ & $-1$ & $0$ & $-1$ \\
Standard$\,\otimes\,$sign ($3$) & $3$ & $-1$ & $-1$ & $0$ & $1$ \\
\bottomrule
\end{tabular}
\end{table}

\begin{table}[h]
\caption{$D_4$ character table (conjugacy classes: $e$, $r^2$, $\{r,r^3\}$,
$\{s,r^2s\}$, $\{rs,r^3s\}$). Verified: $\sum d_i^2=1+1+1+1+4=8=|D_4|$; all
row-orthogonality relations hold exactly.}
\label{tab:d4chars}
\centering
\small
\begin{tabular}{@{}lrrrrr@{}}
\toprule
Irrep (dim) & $e$ & $r^2$ & $\{r,r^3\}$ & $\{s,r^2s\}$ & $\{rs,r^3s\}$ \\
\midrule
Trivial ($1$) & $1$ & $1$ & $1$ & $1$ & $1$ \\
$\chi_2$ ($1$) & $1$ & $1$ & $1$ & $-1$ & $-1$ \\
$\chi_3$ ($1$) & $1$ & $1$ & $-1$ & $1$ & $-1$ \\
$\chi_4$ ($1$) & $1$ & $1$ & $-1$ & $-1$ & $1$ \\
Two-dim.\ ($2$) & $2$ & $-2$ & $0$ & $0$ & $0$ \\
\bottomrule
\end{tabular}
\end{table}

A verified character table is only useful if the projection machinery
built on top of it is also correct. The isotypic-energy decomposition
used throughout \S\ref{sec:h1_results} was additionally verified two
further ways, on synthetic data where the right answer is known in
advance: (i) that the energy fractions it computes sum to exactly $1.0$
on a random embedding, for both groups --- confirming the projections are
a complete decomposition, not one that leaks energy or double-counts it
--- and (ii) that it recovers the exact true dimension when applied to a
synthetic embedding constructed to lie purely within one irrep's isotypic
component, for every irrep of both groups (Table~\ref{tab:s4sanity}). The
second check is the more informative one: it directly tests whether the
decomposition correctly attributes energy to the irrep that actually
generated it, which is exactly what the predicted-bound computation in
\S\ref{sec:h1_results} depends on.

\begin{table}[h]
\caption{$S_4$ isotypic-energy sanity check. A synthetic embedding built
purely from one irrep places $>99\%$ of its energy there and recovers that
irrep's exact dimension as the predicted bound.}
\label{tab:s4sanity}
\centering
\small
\begin{tabular}{@{}lccc@{}}
\toprule
Test & Energy on target irrep & Predicted bound & True dim \\
\midrule
Random embedding & (sums to $1.000$ across all irreps) & -- & -- \\
Pure trivial & $1.000$ & $1$ & $1$ \\
Pure sign & $1.000$ & $1$ & $1$ \\
Pure two-dim.\ & $1.000$ & $2$ & $2$ \\
Pure standard & $1.000$ & $3$ & $3$ \\
\bottomrule
\end{tabular}
\end{table}

\FloatBarrier
\section{Effective Rank Metric Robustness}
\label{app:metric_robustness}

Every result in \S\ref{sec:h1_results} and \S\ref{sec:ablation} depends
on a single choice made early in this project and easy to take for
granted: which of several reasonable definitions of ``effective rank'' to
actually compute at every evaluation step. The naive candidates --- a hard
numerical-rank threshold, or the entropy-based effective rank common in
the representation-geometry literature --- are both far more sensitive to
the small amount of noise present in any real trained network than the
participation-ratio form we use throughout this paper, and using the
wrong one would have made every plateau-rank number in
\S\ref{sec:results} substantially less trustworthy without necessarily
looking wrong on its face. This appendix is the check that motivated the
choice: we inject additive Gaussian noise of increasing scale into a
matrix of known true isotypic dimension $4$ and compare the participation
ratio against the entropy-based effective rank and a hard $99\%$-energy
numerical-rank threshold (Table~\ref{tab:metricrobust}). The contrast is
stark rather than marginal: at a moderate noise level of $0.2$, the
entropy form has already inflated to more than thirteen times the true
value, while the participation ratio has drifted to only about two and a
half times it. The participation ratio degrades far more gracefully across
the entire tested range, which is why it is the statistic reported
everywhere in the main text.

\begin{table}[h]
\caption{Effective rank under additive noise, true isotypic dimension $4$.}
\label{tab:metricrobust}
\centering
\small
\begin{tabular}{@{}lrrrr@{}}
\toprule
Noise scale & Participation ratio & Entropy rank & $99\%$-energy rank & $\sigma_1$ \\
\midrule
$0.00$ & $3.96$ & $3.99$ & $4$ & $23.42$ \\
$0.02$ & $4.01$ & $12.64$ & $4$ & $23.41$ \\
$0.05$ & $4.26$ & $24.80$ & $39$ & $23.42$ \\
$0.10$ & $5.22$ & $38.99$ & $56$ & $23.47$ \\
$0.20$ & $9.65$ & $52.28$ & $61$ & $23.75$ \\
$0.40$ & $28.28$ & $59.87$ & $63$ & $25.06$ \\
\bottomrule
\end{tabular}
\end{table}

\FloatBarrier
\section{Walsh Alignment Score: In-Sample Bias Check}
\label{app:walsh_bias}

H2 (\S\ref{sec:h2}) reports, honestly, that the Walsh alignment score does
not restore a precursor on sparse parity. That negative result is only
trustworthy if the score itself is measuring something real rather than
an artifact of how it is fit, so this appendix documents the specific
methodological problem an earlier version of the score had, and how it
was fixed. The concern is a familiar one in high-dimensional regression:
a degree-$2$ Walsh basis has $821$ features at $n=40$, comfortably enough
to fit noise given a validation set of realistic size, so naively fitting
and scoring the alignment score on the same batch risks reporting a
strong-looking $R^2$ that reflects overfitting rather than genuine
alignment between the hidden representation and the target's spectral
structure. Table~\ref{tab:walshbias} reports the pure-noise $R^2$ for the
in-sample and out-of-sample versions of the score at our configuration's
actual probe sizes ($n=40$ inputs; degree $\le\!1$ has $41$ basis
features, degree $\le\!2$ has $821$). The out-of-sample version --- fit on
half of a dedicated probe batch and scored on the held-out half,
decoupled entirely from the task's own validation set --- returns
$R^2$ indistinguishable from zero on pure noise, as it should. This is
the version used for every Walsh-score number reported anywhere in the
main text.

\begin{table}[h]
\caption{Walsh alignment score bias check, pure-noise inputs.}
\label{tab:walshbias}
\centering
\small
\begin{tabular}{@{}lrr@{}}
\toprule
Configuration & Probe size & Out-of-sample $R^2$ on pure noise \\
\midrule
Degree $\le\!1$ ($n=40$) & $2{,}500$ & $-0.035$ \\
Degree $\le\!2$ ($n=40$) & $20{,}525$ & $-0.087$ \\
\bottomrule
\end{tabular}
\end{table}

The scale of the problem this avoids is worth stating plainly: in-sample
scoring at realistic validation-set sizes for this configuration returned
$R^2$ between $0.46$ and $0.91$ on the same pure-noise inputs across
$k\in\{2,3,4\}$ before this fix --- numbers that would have looked like
strong, genuine alignment on data containing no signal whatsoever. That
gap is the entire reason \S\ref{sec:h2}'s negative result should be
trusted: it was reached with a signal that reliably reports zero on data
containing nothing, not one that would report a false positive under the
slightest encouragement.

\FloatBarrier
\section{Qualitative Failure Diagnostic}
\label{app:s4failure}

The numbers in \S\ref{sec:disentangle} establish that $S_4$ never
converges within budget; this appendix shows what that failure actually
looks like over the course of training, which the aggregate statistics
alone do not convey. Figure~\ref{fig:s4diag} shows accuracy and rank
trajectories for \texttt{s4\_seed0}, the specific run our failure-analysis
logic selects as most informative among every run in the paper that did
not reach the grokking threshold: $S_4$'s failure to converge is
prioritized for this diagnostic over any single non-positive-lead modular
or parity run, since a network stuck at chance accuracy for its entire
training budget is a more informative failure to visualize than a run
that merely lacked a positive precursor lead. The pattern is unambiguous:
train accuracy rises to essentially $1.0$ within the first few hundred
steps, while validation accuracy remains near the $4.2\%$ chance rate for
the full $99{,}999$-step budget. This is ordinary memorization with no
generalization at all --- not a grokking transition that is merely slow
or a run caught mid-transition, but a network that never begins to
generalize in any of the roughly $500$ evaluations logged over its entire
training run.

\begin{figure}[h]
\centering
\includegraphics[width=0.7\textwidth]{figures/figS1_failure_case_diagnostic.pdf}
\caption{Accuracy and participation ratio for \texttt{s4\_seed0} across the
full $100{,}000$-step budget. Train accuracy saturates almost immediately;
validation accuracy never rises above chance.}
\label{fig:s4diag}
\end{figure}


\section{Complete Hyperparameter and Reproducibility Reference}
\label{app:fullconfig}

Every other appendix in this paper documents one specific methodological
decision in depth; this one is the opposite in spirit, and closes the
appendix deliberately for that reason: a single, complete reference table
covering every architecture, optimizer setting, detection threshold,
task configuration, seed, and hardware detail used anywhere in this
paper, in one place, rather than scattered across the sections that each
needed only part of it. Table~\ref{tab:fullconfig} was built by
cross-checking every field against the exact \texttt{config.json} logged
by the run that produced the numbers in \S\ref{sec:results} --- nothing
here is reconstructed from memory or approximated. Values that were swept
rather than held fixed are marked $\star$ and point to the section that
reports the corresponding sweep, so a value that looks unfamiliar in
isolation can always be traced back to the specific experiment that
varied it.

\begin{table}[h]
\caption{Complete hyperparameter and reproducibility reference. $\star$: the
parameter was varied in the sweep named in the Notes column; all unlabeled
values were fixed across every reported run.}
\label{tab:fullconfig}
\centering
\small
\renewcommand{\arraystretch}{1.0}
\begin{tabular}{@{}>{\raggedright\arraybackslash}p{0.27\textwidth}>{\raggedright\arraybackslash}p{0.17\textwidth}>{\raggedright\arraybackslash}p{0.46\textwidth}@{}}
\toprule
\textbf{Parameter} & \textbf{Value} & \textbf{Notes} \\
\midrule
\multicolumn{3}{l}{\textit{Architecture --- modular, $S_4$, $D_4$ (transformer)}} \\
\midrule
Model dimension $d_{\mathrm{model}}$ & $128$ & fixed across all three tasks \\
Transformer layers & $1$ & \\
Attention heads & $4$ & head dim $=32$ \\
FFN dimension $d_{\mathrm{mlp}}$ & $512$ & \\
Input sequence length & $3$ & $[a,b,\texttt{=}]$ or $[g,h,\texttt{=}]$ \\
Positional embedding init & $\mathcal{N}(0,0.02^2)$ & only parameter with a non-default init \\
Other weights (attn.\ / MLP / embed.) & PyTorch default & not separately tuned \\
Parameters & $214{,}656$ / $205{,}184$ / $201{,}088$ & modular ($p{=}113$) / $S_4$ / $D_4$ \\
\midrule
\multicolumn{3}{l}{\textit{Architecture --- sparse parity (MLP)}} \\
\midrule
Hidden width & $200$ & one hidden ReLU layer \\
Parameters & $8{,}602$ & \\
\midrule
\multicolumn{3}{l}{\textit{Optimizer}} \\
\midrule
Algorithm & AdamW & decoupled weight decay \citep{loshchilov2019decoupled} \\
Learning rate $\eta$ & $10^{-3}$ & fixed, all tasks \\
Weight decay $\lambda$ & $2.0$ / $0.7$ & modular, $S_4$, $D_4$ / parity \\
& $\star$ & swept \S\ref{sec:disentangle}: $S_4\in\{0.5,1.0,2.0\}$, $D_4\in\{1.0,2.0\}$ \\
& $\star$ & swept \S\ref{sec:robustness}: modular $\in\{0.5,1.0,3.0,4.0,6.0\}$ \\
Momentum $\beta_1$ & $0.9$ & \\
Second moment $\beta_2$ & $0.98$ & \\
Batch size & full-batch & entire training set per step; varies by task, $38$ ($D_4$) to $5{,}107$ (modular, $p{=}113$) \\
Max.\ training steps & $20$k / $40$k / $100$k & modular / parity / $S_4$, $D_4$ \\
Eval.\ interval & every $200$ steps & \\
\midrule
\multicolumn{3}{l}{\textit{Detection thresholds}} \\
\midrule
Grokking threshold & val.\ acc.\ $\ge 0.95$ & \\
Sustain requirement & $5$ consec.\ evals.\ & plus $15$ additional evals.\ before early stop \\
Clock free parameter $\rho$ & $0.8$ & unfitted default; $\star$ swept \S\ref{sec:clock}: $\{0.1,0.3,0.5,0.8,0.9\}$ \\
Isotypic energy threshold & $90\%$ & fewest irreps whose energy sums to this \\
Fourier top-$k$ (modular) & $4$ & non-DC frequencies \\
Walsh degree bound & $\le\!2$ & degree-$1$ is the primary signal (\S\ref{sec:h2}) \\
\midrule
\multicolumn{3}{l}{\textit{Tasks and data}} \\
\midrule
Modular addition & $(a+b) \bmod p$ & $p\in\{61,97,113\}$; train frac.\ $0.4$ \\
Sparse parity & $\mathrm{sign}(\prod_{i\in S}x_i)$ & $n{=}40$, $k\in\{2,3,4\}$; train frac.\ $0.7$; samples $\{3{,}000,4{,}000,6{,}000\}$ for $k{=}\{2,3,4\}$ \\
$S_4$ composition & order $24$, $576$ pairs & train frac.\ $0.6$ ($345$ train pairs) \\
$D_4$ composition & order $8$, $64$ pairs & train frac.\ $0.6$ ($38$ train pairs) \\
\midrule
\multicolumn{3}{l}{\textit{Seeds and hardware}} \\
\midrule
Seeds & $\{0,\ldots,4\}$ & $5$ per main configuration \\
& $\{0,1,2\}$ & $3$ per robustness-sweep point (\S\ref{sec:robustness}) \\
RNG scope & Python, NumPy, PyTorch & seeded jointly per run \\
Hardware & NVIDIA Tesla T4 & single GPU per run \\
Peak GPU memory & $780$\,MB & \\
Training runs (total) & $70$ & $15$ modular $+$ $15$ parity $+$ $5$ $S_4$ $+$ $5$ $D_4$ $+$ $15$ modular robustness ($5$ wd $\times$ $3$ seeds) $+$ $5$ $D_4$ grid $+$ $10$ $S_4$ grid (new-training seeds only; wd$=2.0$ folded in from the main runs, not retrained) \\
Total wall-clock & $15{,}205$\,s ($4.2$\,h) & mean $217$\,s/run \\
\bottomrule
\end{tabular}
\end{table}


\end{document}
